\subsection{Hilbertian sum of orthogonal subspaces of a hilbertian space}

DEFINITION 2. --- A hilbertian space E is said to be a hilbertian sum of a family \((\mathbf{E}_{i})_{i\in I}\) of closed vector subspaces of E when :

\begin{enumerate}
\def\labelenumi{\arabic{enumi})}
\item
  for two distinct indices \(i\) , \(k\) in I, the subspaces \(\mathrm{E}_i\) and \(\mathrm{E}_k\) are orthogonal in E;
\item
  the closed vector subspace generated by the union of the \(\mathbf{E}_i\) is E.
\end{enumerate}

THEOREM 1. --- Let E be a hilbertian space which is a hilbertian sum of a family \((\mathrm{E}_{i})_{i\in\mathbf{I}}\) of closed vector subspaces of E. There exists an isomorphism f and only one, from E onto the external hilbertian sum \(\bigoplus_{i\in I}E_{i}=F\) of the family \((\mathrm{E}_{i})\) such that, for all \(i\in I\) , the restriction of \(f\) to \(\mathbf{E}\) is the canonical mapping \(f_{i}\) from \(\mathbf{E}_i\) into \(\mathbf{F}\) .

Let \(S \subset F\) be the « algebraic » direct sum of the \(E_i\) , and let \(g\) be the linear mapping \((x_i)_{i \in I} \mapsto \sum_{i \in I} x_i\) from \(S\) into \(E\) . We shall show that \(g\) is an isomorphism from the prehilbertian space S onto the (prehilbertian) subspace \(g(S)\) of E, generated by the union of the \(E_{i}\) : for, for two elements \(\mathbf{x} = (x_{i})_{i \in \mathrm{I}}\) , \(\mathbf{y} = (y_{i})_{i \in \mathrm{I}}\) , we have

\[
\langle g (\mathbf {x}) | g (\mathbf {y}) \rangle = \left\langle \sum_ {i \in \mathrm{I}} x _ {i} \right| \sum_ {i \in \mathrm{I}} y _ {i} \rangle = \sum_ {(i, k) \in \mathrm{I} \times \mathrm{I}} \left\langle x _ {i} \mid y _ {k} \right\rangle .
\]

But if \(i \neq k\) , \(\langle x_i | y_k \rangle = 0\) by hypothesis, hence

\[
\langle g (\mathbf {x}) | g (\mathbf {y}) \rangle = \sum_ {i \in \mathbf {I}} \langle x _ {i} | y _ {i} \rangle = \langle \mathbf {x} | \mathbf {y} \rangle ;
\]

this proves our assertion. Since S is dense in F and \(g(S)\) dense in E, the isomorphism g extends to an isomorphism \(\overline{g}\) from F onto E (V, p.~8, cor.). It is clear that the inverse isomorphism f of \(\overline{g}\) is the required mapping; its uniqueness follows from the fact that the closed subspace of E generated by the union of the \(E_{i}\) is E itself.

When E is the hilbertian sum of a family \((\mathrm{E}_{i})_{i\in I}\) of subspaces, we shall often identify E with the external hilbertian sum F of the \(E_{i}\) by means of the isomorphism f.~If the set I is finite, saying that E is the hilbertian sum of the family \((\mathrm{E}_{i})_{i\in I}\) means that the \(E_{i}\) are two by two orthogonal and that the vector space E is the direct sum of the family \((\mathrm{E}_{i})_{i\in I}\) of subspaces.

COROLLARY 1. --- Let E be a hilbertian space, which is a hilbertian sum of a family \((\mathbf{E}_{i})_{i\in\mathbf{I}}\) of closed vector subspaces of E; for all \(i\in I\) , let \(p_{E_{i}}\) be the orthoprojector (V, p.~13) from E onto \(E_{i}\) .

\begin{enumerate}
\def\labelenumi{\alph{enumi})}
\tightlist
\item
  For all \(x \in \mathbb{E}\) , the family \((\| p_{\mathbf{E}_i}(x)\|^{2})_{i\in \mathbf{I}}\) is summable in \(\mathbf{R}\) , the family \((p_{\mathbf{E}_i}(x))_{i\in \mathbf{I}}\) is summable in \(\mathbb{E}\) , and we have
\end{enumerate}

\[
\left\| x \right\| ^ {2} = \sum_ {i \in \mathbf {I}} \left\| p _ {\mathrm{E} _ {i}} (x) \right\| ^ {2}, \quad x = \sum_ {i \in \mathbf {I}} p _ {\mathrm{E} _ {i}} (x).
\]

\begin{enumerate}
\def\labelenumi{\alph{enumi})}
\setcounter{enumi}{1}
\item
  Conversely, if \((x_{i})_{i\in \mathbf{I}}\) is a family of elements of \(\mathbf{E}\) such that \(x_{i}\in E_{i}\) for all \(i\in \mathbf{I}\) and \(\sum_{i\in \mathbf{I}}\| x_i\|^2 < + \infty\) , this family is summable, and the sum \(x\) is the only point of \(\mathbf{E}\) for which \(p_{\mathrm{E}_i}(x) = x_i\) for all \(i\in \mathbf{I}\) .
\item
  For every pair of points \(x, y\) of \(\mathbf{E}\) , we have
\end{enumerate}

\[
\langle x | y \rangle = \sum_ {i \in \mathbf {I}} \left\langle p _ {\mathrm{E} _ {i}} (x) \mid p _ {\mathrm{E} _ {i}} (y) \right\rangle .
\]

These properties are in fact obvious for the external hilbertian sum of the \(\mathbf{E}_i\) , and can be transferred to \(\mathbf{E}\) by isomorphism.

COROLLARY 2. --- Let E be a Hausdorff prehilbertian space, \((\mathrm{E}_{i})_{i\in\mathrm{I}}\) a family of complete vector subspaces of E such that, for every pair of distinct indices i, k in I, the subspaces \(E_{i}\) and \(E_{k}\) are orthogonal. Let V be the closed vector subspace of E generated by the union of the \(E_{i}\) . For every \(i\in I\) , let \(p_{E_{i}}\) be the orthoprojector from E onto \(E_{i}\) . Let \(x\in E\) .

\begin{enumerate}
\def\labelenumi{\arabic{enumi})}
\item
  We have \(\sum\left\|p_{E_{i}}(x)\right\|^{2}\leqslant\left\|x\right\|^{2}\) .
\item
  The following conditions are equivalent: a) \(x \in \mathbf{V}\) ; b) \(\sum_{i \in I} \| p_{\mathrm{E}_i}(x) \|^2 = \| x \|^2\) ; c) the family \((p_{\mathrm{E}_i}(x))_{i \in I}\) is summable in E, and we have \(x = \sum_{i \in I} p_{\mathrm{E}_i}(x)\) .
\item
  Suppose V is complete. Then the family \((p_{\mathrm{E}_i}(x))_{i\in \mathbf{I}}\) is summable in E, and
\end{enumerate}

\[
p _ {\mathbf {V}} (x) = \sum_ {i \in \mathbf {I}} p _ {\mathrm{E} _ {i}} (x), \quad \left\| p _ {\mathbf {V}} (x) \right\| ^ {2} = \sum_ {i \in \mathbf {I}} \left\| p _ {\mathrm{E} _ {i}} (x) \right\| ^ {2},
\]

where \(p_{V}\) denotes the orthoprojector from E onto V.

Let \(\hat{\mathbf{E}}\) be the hilbertian space completion of \(\mathbf{E}\) ; we identify \(\mathbf{E}\) with a dense subspace of \(\hat{\mathbf{E}}\) ; the \(\mathbf{E}_i\) , being complete, are closed subspaces of \(\hat{\mathbf{E}}\) . The closure \(\overline{\mathbf{V}}\) of \(\mathbf{V}\) in \(\hat{\mathbf{E}}\) is the closed vector subspace of \(\hat{\mathbf{E}}\) generated by the union of the \(\mathbf{E}_i\) , and \(\mathbf{V} = \overline{\mathbf{V}} \cap \mathbf{E}\) . The space \(\hat{\mathbf{E}}\) is the hilbertian sum of the \(\mathbf{E}_i\) and of the subspace \(\mathbf{W}\) , the orthogonal complement of \(\overline{\mathbf{V}}\) in \(\hat{\mathbf{E}}\) ; put \(x_0 = p_{\mathrm{W}}(x)\) and \(x_i = p_{\mathrm{E}_i}(x)\) for all \(i \in \mathbf{I}\) . By cor. 1, we have \(\| x\|^2 = \| x_0\|^2 + \sum_{i \in \mathbf{I}} \| x_i\|^2\) , and \(x = x_0 + \sum_{i \in \mathbf{I}} x_i\) in \(\hat{\mathbf{E}}\) . This implies assertion 1), and the fact that conditions \(b)\) and \(c)\) of 2) are equivalent to the condition \(x_0 = 0\) , hence to the condition \(x \in \mathbf{V}\) . Finally, if \(\mathbf{V}\) is complete, and if we put \(x' = p_{\mathrm{V}}(x)\) , we have \(x' - x_i = (x - x_i) - (x - p_{\mathrm{V}}(x))\) , hence \(x' - x_i\) is orthogonal to \(\mathbf{E}_i\) , and so \(x_i = p_{\mathrm{E}_i}(x')\) for all \(i \in \mathbf{I}\) ; it is now enough to apply property 2) to the vector \(x'\) .

Remark. --- Let E be a Hausdorff prehilbertian space, \((\mathrm{V}_{i})_{i\in\mathrm{I}}\) a family of vector subspaces of E such that for every pair of distinct indices i, k, the subspaces \(V_{i}\) and \(V_{k}\) are orthogonal. Then, for every \(k\in I\) , the intersection of \(V_{k}\) and of the closed vector subspace \(W_{k}\) generated by the union of the \(V_{i}\) for all \(i\neq k\) reduces to 0 for, if x belongs to \(V_{k}\) and also to \(W_{k}\) , then it is orthogonal to all the \(V_{i}\) for \(i\neq k\) , hence to \(W_{k}\) . In particular, x is orthogonal to itself, hence is zero.

PROPOSITION 2. --- Let E be a hilbertian space and \((\mathrm{V}_{\lambda})_{\lambda\in\mathrm{L}}\) a family of closed vector subspaces of E; for every \(\lambda\in L\) , let \((\mathrm{W}_{\lambda\mu})_{\mu\in\mathrm{M}_{\lambda}}\) be a family of closed vector subspaces of \(V_{\lambda}\) such that \(V_{\lambda}\) is the closed vector subspace generated by the union of this family. In order that E is the hilbertian sum of the family \((\mathrm{W}_{\lambda\mu})_{\lambda\in\mathrm{L},\mu\in\mathrm{M}_{\lambda}}\) , it is necessary and sufficient that E is the hilbertian sum of the family \((\mathrm{V}_{\lambda})_{\lambda\in\mathrm{L}}\) and that, for each \(\lambda\in L\) , \(V_{\lambda}\) is the hilbertian sum of the family \((\mathrm{W}_{\lambda\mu})_{\mu\in\mathrm{M}_{\lambda}}\) (« associativity of the hilbertian sum »). To show that the condition is necessary, it is enough to see that \(V_{\alpha}\) and \(V_{\beta}\) are orthogonal if \(\alpha\neq\beta\) . But, every element of \(W_{\alpha\mu}\) ( \(\mu\in M_{\alpha}\) ) is orthogonal to all the \(W_{\beta\nu}\) ( \(\nu\in M_{\beta}\) ), hence to the closed vector subspace \(V_{\beta}\) which they generate; the same argument then shows that every element of \(V_{\beta}\) is orthogonal to \(V_{\alpha}\) , being orthogonal to all the \(W_{\alpha\mu}\) ( \(\mu\in M_{\alpha}\) ).

To show that the condition is sufficient, it is enough to verify that, if it is satisfied, \(\mathrm{E}\) is equal to the closed vector subspace \(\mathrm{F}\) generated by the union of the \(\mathbf{W}_{\lambda \mu}\) (\(\lambda \in \mathbf{L}\), \(\mu\in M_{\lambda}\)); but, for each \(\lambda\in L\) , F contains the closed vector subspace generated by the union of the \(W_{\lambda\mu}\) such that \(\mu\in M_{\lambda}\) , that is, F contains \(V_{\lambda}\) ; hence F is the closed vector subspace generated by the union of the \(V_{\lambda}\) , which is E by hypothesis.

